%% toc: P3 27 27.1,27.2,27.3,27.4
\section{Horn formulas and the construction behind isomorphic ultrapowers}
\label{sec:p03-w03-p03-017}

This section supplies two arguments that the surrounding treatment of
products and ultraproducts states but does not carry out.  The first is the
preservation of Horn formulas under direct and reduced products, with the
equality and false-head conventions made explicit.  The second is the
complete construction, without any continuum hypothesis, of the single
ultrafilter in Shelah's isomorphic-ultrapowers theorem
(\cref{thm:p03-w03-p03-018-shelah}), together with written proofs of both
directions of \cref{thm:p03-keisler-shelah}.  The last two subsections add
the regularity of Shelah's bound, the earlier approaches that Shelah
describes, and those of Keisler's 1961 characterizations of elementary
classes that his theorem frees from the continuum hypothesis.

\emph{Inputs and conclusions.}  The inputs are first-order structures of a
common signature, a nonempty index set with a proper filter on it, and, for
the second part, an infinite cardinal \(\lambda\).  The conclusions are
truth of Horn formulas in reduced products, an ultrafilter \(D\) over
\(\lambda\) with isomorphic ultrapowers for all elementarily equivalent pairs
of small power, saturation of those ultrapowers, and the Keisler--Shelah
equivalence.  The ambient set theory is \(\mathsf{ZFC}\); the uses of Choice
are those recorded in \cref{def:p03-ultrafilter-choice-boundary}, and each
proof below names the step at which Choice enters.

\subsection{Horn formulas and preservation under products}
\label{sec:p03-w03-p03-017-horn}

The syntax follows \textcite{BurrisSankappanavar2012}.  Its atomic formulas
of type \(L\) over a set of variables are the equations \(p\approx q\)
between terms and the relational atoms \(r(p_1,\ldots,p_n)\)
\cite[V\S1, Definition~1.4, p.~192]{BurrisSankappanavar2012}.  Equality is
therefore not a relation symbol of the signature: it is the logical symbol
\(\approx\), interpreted in every structure as identity.

\begin{hbdefinition}{def:p03-w03-p03-017-horn}{Basic Horn formulas and Horn formulas}
A \emph{literal} is an atomic formula or the negation of one; equations
\(p\approx q\) are atomic.  A \emph{basic Horn formula} is a disjunction
\(\beta_1\vee\cdots\vee\beta_k\), with \(k\ge1\), of literals of which at
most one is atomic.  A \emph{Horn formula} is a formula in prenex form
\[
  Q_1y_1\cdots Q_my_m\,(\varphi_1\wedge\cdots\wedge\varphi_r),
  \qquad r\ge1,\ m\ge0,
\]
whose matrix is a conjunction of basic Horn formulas and whose
\(Q_j\) are quantifiers
\cite[V\S2, Definition~2.5, p.~207]{BurrisSankappanavar2012}; the same
definition is given before Exercise~29 of
\cite[V\S1, p.~204]{BurrisSankappanavar2012}.  A basic Horn formula with
exactly one atomic disjunct \(\alpha\) is written
\(\alpha_1\wedge\cdots\wedge\alpha_{k-1}\to\alpha\); one with no atomic
disjunct is written \(\neg(\alpha_1\wedge\cdots\wedge\alpha_k)\) and is said
to have a \emph{false head}.
\index{Horn formula}\index{basic Horn formula}
\end{hbdefinition}

\begin{hbconvention}{def:p03-w03-p03-017-horn-conventions}{Equality, false heads and the constants \(\bot\), \(\top\)}
\begin{enumerate}
  \item \emph{Equality.}  An equation \(p\approx q\) may occur as the
        atomic disjunct of a basic Horn formula or negated among its other
        disjuncts.  Thus the cancellation law
        \(\forall x\forall y\forall z\,(x\cdot z\approx y\cdot z\to x\approx
        y)\) is a Horn sentence; expressing this law by Horn formulas is
        the first item of
        \cite[V\S1, Exercise~29, p.~204]{BurrisSankappanavar2012}.  In a
        reduced product the symbol \(\approx\) denotes identity of classes,
        which is the filter congruence of
        \cref{def:p03-w03-p03-017-reduced}; no separate equality axiom is
        imposed.
  \item \emph{False head.}  A disjunction of negated atomic formulas alone
        is a basic Horn formula; it is how a Horn formula forbids a
        configuration.  The
        empty disjunction (\(k=0\)) is not a basic Horn formula.
  \item \emph{The constants \(\bot\) and \(\top\).}  The source's syntax has
        no propositional constants.  The class \(\mathrm{HORN}\) of
        \cref{def:p02-w02b-p02-031-krom-horn} writes the false head as an
        implication with conclusion \(\mathit{false}\); that clause is read
        here as the basic Horn formula with no atomic disjunct.  If a
        language does admit \(\bot\) and \(\top\) as atomic formulas, their
        truth sets are \(\varnothing\) and \(I\), and the atomic condition of
        \cref{thm:p03-w03-p03-017-atomic} holds for them as well; every
        result below is unchanged.
  \item \emph{Comparison with the decision-problem class.}  The class
        \(\mathrm{HORN}\) of \cref{def:p02-w02b-p02-031-krom-horn} is
        stated there for languages without equality and function symbols.
        On such languages it coincides with the Horn formulas defined here;
        the present definition also admits function symbols and equations.
\end{enumerate}
\end{hbconvention}

\begin{hbdefinition}{def:p03-w03-p03-017-reduced}{Direct and reduced products}
Let \(\langle A_i:i\in I\rangle\) be a \emph{nonempty} indexed family of
structures of type \(L\), each with nonempty universe.  The \emph{direct
product} \(\prod_{i\in I}A_i\) has universe the Cartesian product, with
operations computed coordinatewise and with \(r(a_1,\ldots,a_n)\) holding
when \(r(a_1(i),\ldots,a_n(i))\) holds in \(A_i\) for every \(i\in I\)
\cite[V\S1, p.~204]{BurrisSankappanavar2012}.

Let \(F\) be a \emph{proper} filter over \(I\) (so \(\varnothing\notin F\)).
For tuples from the product and a formula \(\varphi\) write
\[
  [\![\varphi(a_1,\ldots,a_n)]\!]
  =\{i\in I: A_i\models\varphi(a_1(i),\ldots,a_n(i))\},
\]
the local truth set of \cref{def:p03-fv-product}.  Put
\(\langle a,b\rangle\in\theta_F\) iff \([\![a\approx b]\!]\in F\).  The
\emph{reduced product} \(\prod_{i\in I}A_i/F\) has universe the set of
classes \(a/F\) of \(\theta_F\); a function symbol \(f\) is interpreted by
\(f(a_1/F,\ldots,a_n/F)=f(a_1,\ldots,a_n)/F\), and a relation symbol \(r\)
by: \(r(a_1/F,\ldots,a_n/F)\) holds iff
\([\![r(a_1,\ldots,a_n)]\!]\in F\)
\cite[V\S2, Definitions~2.1 and~2.3, pp.~205--206]{BurrisSankappanavar2012}.
When \(F\) is an ultrafilter this is the ultraproduct of
\cref{sec:types-ultraproducts-05}.
\index{reduced product}\index{direct product!of structures}
\end{hbdefinition}

The filter \(F\) is the dual of the proper ideal \(\{I\setminus S:S\in F\}\),
and \(\prod_{i\in I}A_i/F\) is the structure that
\cref{def:p03-fv-reduced-standing} writes with that ideal.  The direct
product is the reduced product by the filter \(\{I\}\), which is proper
because \(I\neq\varnothing\): then \(\theta_{\{I\}}\) is identity, and
\(a\mapsto a/\{I\}\) is an isomorphism
(\cref{thm:p03-w03-p03-017-atomic}).

\begin{hblemma}{thm:p03-w03-p03-017-atomic}{proved=annex:p03-w03-p03-017-atomic}{The congruence and the atomic condition}
With the data of \cref{def:p03-w03-p03-017-reduced}:
\begin{enumerate}
  \item \(\theta_F\) is an equivalence relation on \(\prod_{i\in I}A_i\),
        compatible with every fundamental operation, so the interpretations
        of \cref{def:p03-w03-p03-017-reduced} do not depend on the
        representatives chosen;
  \item for every term \(t(\bar x)\),
        \(t^{\prod A_i/F}(\bar a/F)=b_t/F\) where
        \(b_t(i)=t^{A_i}(\bar a(i))\);
  \item for every atomic formula \(\alpha(\bar x)\), equations included,
        \(\prod_{i\in I}A_i/F\models\alpha(\bar a/F)\) iff
        \([\![\alpha(\bar a)]\!]\in F\);
  \item \(a\mapsto a/\{I\}\) is an isomorphism of \(\prod_{i\in I}A_i\) onto
        \(\prod_{i\in I}A_i/\{I\}\).
\end{enumerate}
\end{hblemma}

\begin{hblemma}{thm:p03-w03-p03-017-maximal}{proved=annex:p03-w03-p03-017-maximal}{The maximal property}
Assume the Axiom of Choice.  Let \(\langle A_i:i\in I\rangle\) be a nonempty
family of structures with nonempty universes, let
\(\varphi(x,y_1,\ldots,y_n)\) be a formula and let
\(a_1,\ldots,a_n\in\prod_iA_i\).  There is \(a\in\prod_iA_i\) with
\([\![\exists x\,\varphi(x,\bar a)]\!]=[\![\varphi(a,\bar a)]\!]\)
\cite[V\S2, Lemma~2.6, p.~207]{BurrisSankappanavar2012}.
\end{hblemma}

\begin{hbtheorem}{thm:p03-w03-p03-017-horn-preservation}{proved=annex:p03-w03-p03-017-horn-preservation}{Horn formulas are preserved under reduced products}
Assume the Axiom of Choice.  Let \(\prod_{i\in I}A_i/F\) be a reduced
product of a nonempty family of structures of type \(L\) by a proper filter
\(F\), let \(\varphi(x_1,\ldots,x_n)\) be a Horn formula of type \(L\), and
let \(a_1,\ldots,a_n\in\prod_iA_i\).  If
\([\![\varphi(a_1,\ldots,a_n)]\!]\in F\), then
\(\prod_{i\in I}A_i/F\models\varphi(a_1/F,\ldots,a_n/F)\)
\cite[V\S2, Theorem~2.7, pp.~207--210]{BurrisSankappanavar2012}.
\end{hbtheorem}

\begin{hbcorollary}{thm:p03-w03-p03-017-horn-products}{proved=annex:p03-w03-p03-017-horn-products}{Horn sentences in direct products, reduced products and ultraproducts}
Assume the Axiom of Choice.  Let \(\sigma\) be a Horn sentence and
\(\langle A_i:i\in I\rangle\) a nonempty family of structures.
\begin{enumerate}
  \item If \(\{i:A_i\models\sigma\}\in F\) for a proper filter \(F\) over
        \(I\), then \(\prod_iA_i/F\models\sigma\).  In particular a Horn
        sentence true in every factor is true in every reduced product and
        every ultraproduct of the family.
  \item If \(A_i\models\sigma\) for every \(i\in I\), then
        \(\prod_{i\in I}A_i\models\sigma\); this is
        \cite[V\S1, Exercise~30, p.~204]{BurrisSankappanavar2012}.
  \item A sentence logically equivalent to a Horn sentence has the same
        preservation properties.
\end{enumerate}
\end{hbcorollary}

\begin{hbexample}{def:p03-w03-p03-017-horn-example}{A Horn sentence, a non-Horn sentence, and a direct product}
Let \(L\) have two unary relation symbols \(P\), \(Q\).  The sentence
\(\forall x\,\neg(P(x)\wedge Q(x))\) is Horn with a false head.  The sentence
\(\sigma=\forall x\,(P(x)\vee Q(x))\) is not Horn: its matrix has two atomic
disjuncts.  Let \(A_1\) have universe \(\{0\}\) with \(P\) true and \(Q\)
false, and \(A_2\) universe \(\{0\}\) with \(P\) false and \(Q\) true.  Both
factors satisfy \(\sigma\).  In \(A_1\times A_2\) the only element is
\((0,0)\); \(P(0,0)\) fails because it fails in the second coordinate, and
\(Q(0,0)\) fails because it fails in the first.  So
\(A_1\times A_2\not\models\sigma\).  The false-head sentence holds in both
factors and, by \cref{thm:p03-w03-p03-017-horn-products}, in the product:
indeed neither relation holds at \((0,0)\).  For the equational case, the
cancellation law of \cref{def:p03-w03-p03-017-horn-conventions} holds in the
additive monoid \(\mathbb N\) and therefore in every reduced power of it.
\end{hbexample}

\begin{hbwarning}{def:p03-w03-p03-017-horn-limits}{Where the hypotheses are needed}
\begin{itemize}
  \item \emph{Proper filter.}  For the improper filter \(\mathcal P(I)\) all
        elements are identified and every atomic formula holds in the
        one-element quotient.  The Horn sentence
        \(\exists x\,\neg R(x)\) can hold in every factor and fail there.
        This is why \(F\) is assumed proper
        \cite[V\S2, Definition~2.1, p.~205]{BurrisSankappanavar2012}.
  \item \emph{Nonempty family.}  For \(I=\varnothing\) the only filter is
        \(\{\varnothing\}\), which is improper, and the empty direct product
        is a one-element structure in which every atomic formula holds.  The
        same sentence \(\exists x\,\neg R(x)\) is then vacuously true in all
        factors and false in the product.
  \item \emph{Choice.}  The existential step uses the maximal property,
        that is, a choice of witnesses across the index set.
  \item \emph{Converse.}  The result is one direction only.  The theorem
        that a sentence preserved under reduced products is equivalent to a
        Horn sentence is due to Keisler, as Shelah reports
        \cite[p.~225]{Shelah1971Ultrapowers}; it is not proved here.
\end{itemize}
\end{hbwarning}

\paragraph{Relationships.}  Atomic satisfaction in direct products is
\cref{thm:p03-fv-coordinate-terms}; arbitrary first-order formulas in
reduced products are handled by the Feferman--Vaught reduction of
\cref{thm:p03-fv-reduced-reduction}, of which the present theorem is the
special case that needs no Boolean condition beyond membership in \(F\).
For ultrafilters, \cref{thm:p03-w03-p03-017-horn-preservation} is one half
of \L o\'s's theorem (\cref{thm:los}) restricted to Horn formulas; for proper
filters that are not ultrafilters \L o\'s's theorem fails, as the example of
\cref{thm:p03-fv-reduced-example} shows.  Subdirect products and their
equational and quasi-equational classes are treated in
\cref{def:p10-subdirect-product}.

\subsection{The construction of an ultrafilter for isomorphic ultrapowers}
\label{sec:p03-w03-p03-017-shelah}

This subsection presents the proof in
\textcite{Shelah1971Ultrapowers} of parts~1 and~2 of
\cref{thm:p03-w03-p03-018-shelah}.  The cardinals \(\lambda\) and
\(\mu=\min\{\mu:\lambda^\mu>\lambda\}\) and the notation \(\dom M\),
\(\lVert M\rVert\), \(M^\lambda/D\) are those of
\cref{def:p03-w03-p03-018-notation}; recall
(\cref{thm:p03-w03-p03-018-mu}) that \(\aleph_0\le\mu\le\lambda\) and
\(\lambda^\nu=\lambda\) for \(1\le\nu<\mu\).  The proof uses only the
definition of an ultrapower and \L o\'s's theorem
\cite[p.~225]{Shelah1971Ultrapowers}.  It generalizes Kunen's elimination of
the generalized continuum hypothesis from the existence of good
ultrafilters, as the source says \cite[pp.~224--225]{Shelah1971Ultrapowers}.

The notation paragraph of the source adds, without proof, that \(\mu\) is a
regular cardinal \cite[p.~225]{Shelah1971Ultrapowers}.  The construction
below does not use this; it is recorded and proved here because it is part
of the source's description of \(\mu\).

\begin{hblemma}{thm:p03-w03-p03-017-mu-regular}{proved=annex:p03-w03-p03-017-mu-regular}{The bound \(\mu\) is regular}
Assume the Axiom of Choice.  Let \(\lambda\) be an infinite cardinal and
\(\mu=\min\{\mu:\lambda^\mu>\lambda\}\).  Then \(\mu\) is a regular cardinal
in the sense of \cref{def:p06-cofinality-regular-singular}
\cite[p.~225]{Shelah1971Ultrapowers}.
\end{hblemma}

\begin{hbconvention}{def:p03-w03-p03-017-shelah-notation}{Further notation of the source}
Following \cite[p.~225]{Shelah1971Ultrapowers}: \(\chi\), \(\kappa\) denote
cardinals, and in \cref{def:p03-w03-p03-017-consistent} and the lemmas
that use it \(\kappa\) is infinite; \(i,j,k,l,\alpha,\beta,\gamma,\delta\)
denote ordinals; \(m,n\) natural numbers.  The letter \(f\) ranges over functions from \(\lambda\) into
\(\mu\), and \(g\) over functions from \(\lambda\) into some cardinal
\(\chi(g)<\mu\); \(F\), \(G\) are sets of such functions, and functions
\(f\in F\) carrying different indices are different functions.  \(D\) is a
proper filter over \(\lambda\).  For a family \(E\) of subsets of
\(\lambda\),
\[
  [E]=\Bigl\{A\subseteq\lambda:\text{for some }A_1,\ldots,A_n\in E,\
  \bigcap_{m=1}^nA_m\subseteq A\Bigr\}
\]
is the filter generated by \(E\), with the empty intersection read as
\(\lambda\).  \(A=\varnothing\ (\mathrm{mod}\ D)\) means that
\(A\subseteq\lambda-X\) for some \(X\in D\); so
\(A\neq\varnothing\ (\mathrm{mod}\ D)\) means that \(A\) meets every member
of \(D\).  \((\dom M)^\lambda\), which the source writes with the universe
between single bars, is the set of functions from \(\lambda\) into
\(\dom M\); for \(\bar a\in(\dom M)^\lambda\) and \(k<\lambda\), \(\bar a[k]\)
is its value at \(k\).
\end{hbconvention}

\begin{hbdefinition}{def:p03-w03-p03-017-consistent}{\(\kappa\)-consistency}
A triple \((F,G,D)\) is \emph{\(\kappa\)-consistent} if
\begin{enumerate}
  \item[A)] \(D\) is generated by a family of at most \(\kappa\) subsets of
        \(\lambda\); and
  \item[B)] whenever \(f^i\in F\) and \(j_i<\mu\) for \(i<\chi<\mu\), and
        \(f^m\in F\), \(g^m\in G\) for \(m\le n\), with all the listed
        \(f\)'s distinct, the set of \(k<\lambda\) with \(f^i(k)=j_i\) for
        all \(i<\chi\) and \(f^m(k)=g^m(k)\) for all \(m\le n\) is
        \(\neq\varnothing\ (\mathrm{mod}\ D)\).
\end{enumerate}
\cite[p.~226, Definition~1]{Shelah1971Ultrapowers}.  The case \(\chi=0\)
with no pairs \((f^m,g^m)\) is included; there the displayed set is
\(\lambda\), and B) says that \(D\) is proper
\cite[p.~230]{Shelah1971Ultrapowers}.
\end{hbdefinition}

\begin{hblemma}{thm:p03-w03-p03-017-lemma1}{proved=annex:p03-w03-p03-017-lemma1}{Lemma~1: an independent family}
There is a family \(F\) of \(2^\lambda\) functions from \(\lambda\) to
\(\mu\) such that \((F,\varnothing,\{\lambda\})\) is \(\mu\)-consistent
\cite[p.~226, Lemma~1]{Shelah1971Ultrapowers}.  The source notes that this
generalizes Ketonen's lemma used by Kunen and that both already appear in
Engelking and Kar\l owicz.
\end{hblemma}

\begin{hblemma}{thm:p03-w03-p03-017-lemma2}{proved=annex:p03-w03-p03-017-lemma2}{Lemma~2: monotonicity and limits}
\begin{enumerate}
  \item[A)] If \((F,G,D)\) is \(\kappa\)-consistent and
        \(\kappa\le\kappa_1\), then \((F,G,D)\) is
        \(\kappa_1\)-consistent.
  \item[B)] If \((F_i,G_i,D_i)\) is \(\kappa_i\)-consistent for every
        \(i<\delta\); \(D_i\subseteq D_j\), \(G_i\subseteq G_j\) and
        \(F_i\supseteq F_j\) for \(i<j<\delta\); \(D=\bigcup_{i<\delta}D_i\),
        \(G=\bigcup_{i<\delta}G_i\), \(F=\bigcap_{i<\delta}F_i\); and
        \(\kappa\ge\kappa_i\) for every \(i<\delta\) and
        \(\kappa\ge\mathrm{cf}(\delta)\), then \((F,G,D)\) is
        \(\kappa\)-consistent.
  \item[C)] If \((F,G,D)\) is \(\kappa\)-consistent, \(F'\subseteq F\) and
        \(G'\subseteq G\), then \((F',G',D)\) is \(\kappa\)-consistent.
\end{enumerate}
\cite[pp.~226--227, Lemma~2]{Shelah1971Ultrapowers}.  In B), \(\delta\) is a
limit ordinal.
\end{hblemma}

\begin{hblemma}{thm:p03-w03-p03-017-lemma3}{proved=annex:p03-w03-p03-017-lemma3}{Lemma~3: bounded-range functions}
Suppose that \((F,\varnothing,D)\) is \(\kappa\)-consistent, that
\(\mu+\abs G\le\kappa\), and that \(G\) is a set of functions from \(\lambda\) to cardinals \({<}\,\mu\).
Then there is \(F'\subseteq F\) with \(\abs{F-F'}\le\kappa\) such that
\((F',G,D)\) is \(\kappa\)-consistent
\cite[p.~227, Lemma~3]{Shelah1971Ultrapowers}.
\end{hblemma}

\begin{hblemma}{thm:p03-w03-p03-017-lemma4}{proved=annex:p03-w03-p03-017-lemma4}{Lemma~4: deciding sets}
\begin{enumerate}
  \item[A)] Suppose \((F,G,D)\) is \(\kappa\)-consistent and
        \(A\subseteq\lambda\).  Then there is \(F'\subseteq F\) with
        \(\abs{F-F'}<\mu\) such that \((F',G,[D\cup\{A\}])\) is
        \(\kappa\)-consistent or \((F',G,[D\cup\{\lambda-A\}])\) is
        \(\kappa\)-consistent.
  \item[B)] If \((F,G,D)\) is \(\kappa\)-consistent, \(A_\alpha\subseteq\lambda\)
        for \(\alpha<\kappa\), and \(\mu\le\kappa\), then there are
        \(F'\subseteq F\) with \(\abs{F-F'}\le\kappa\) and a filter
        \(D'\supseteq D\) such that \((F',G,D')\) is \(\kappa\)-consistent
        and, for every \(\alpha<\kappa\), either \(A_\alpha\in D'\) or
        \(\lambda-A_\alpha\in D'\).
\end{enumerate}
\cite[p.~228, Lemma~4]{Shelah1971Ultrapowers}.
\end{hblemma}

\begin{hblemma}{thm:p03-w03-p03-017-lemma5}{proved=annex:p03-w03-p03-017-lemma5}{Lemma~5: realizing a type by a new function}
Let \(M\) be a model of cardinality \(\chi<\mu\), let
\(\bar a_{l,m}\in(\dom M)^\lambda\) for \(l<l_0\) and \(1\le m\le n_l\), let
\(\mu\le\kappa\), and let \((F,\varnothing,D)\) be \(\kappa\)-consistent with
\(\abs F>\kappa\).  Let
\(p=\{\phi_l(x,y_{l,1},\ldots,y_{l,n_l}):l<l_0<\kappa^+\}\) be formulas of
the language of \(M\), closed under conjunctions in the sense that for
\(l,l'<l_0\) some \(l''<l_0\) has
\(M\models\phi_{l''}[c,\bar a_{l''}[k]]\leftrightarrow
(\phi_l[c,\bar a_l[k]]\wedge\phi_{l'}[c,\bar a_{l'}[k]])\) for all
\(c\in\dom M\) and \(k<\lambda\), where \(\bar a_l[k]\) abbreviates
\(\bar a_{l,1}[k],\ldots,\bar a_{l,n_l}[k]\).  Suppose that for every
\(l<l_0\)
\[
  A^l=\{k<\lambda:M\models(\exists x)\phi_l(x,\bar a_l[k])\}\in D.
\]
Then there are \(\bar a\in(\dom M)^\lambda\), \(F'\subseteq F\) and a filter
\(D'\supseteq D\) such that \(\abs{F-F'}\le\kappa\), \((F',\varnothing,D')\)
is \(\kappa\)-consistent, and
\(\{k<\lambda:M\models\phi_l[\bar a[k],\bar a_l[k]]\}\in D'\) for every
\(l<l_0\) \cite[p.~229, Lemma~5]{Shelah1971Ultrapowers}.
\end{hblemma}

\begin{hblemma}{thm:p03-w03-p03-017-language}{proved=annex:p03-w03-p03-017-language}{Reduction to a fixed language of cardinality \(\lambda\)}
Fix a language \(L^0\) with \(\lambda\) constant symbols and, for each
\(n\ge1\), \(\lambda\) relation symbols and \(\lambda\) function symbols of
arity \(n\).  Let \(M\equiv N\) be models of a language \(L\) with
\(\lVert M\rVert,\lVert N\rVert<\mu\).  There are a language
\(L_*\subseteq L^0\) and \(L_*\)-models \(M_*\equiv N_*\) whose universes are
the cardinals \(\lVert M\rVert\) and \(\lVert N\rVert\), with the following
property: for every ultrafilter \(D\) over \(\lambda\), if
\(M_*^\lambda/D\cong N_*^\lambda/D\) then \(M^\lambda/D\cong N^\lambda/D\),
and if every set of formulas in finitely many variables with at most
\(\kappa\) parameters that is finitely satisfiable in \(M_*^\lambda/D\) is
realized there, then the same holds for \(M^\lambda/D\).  The number of pairs \((M_*,N_*)\) of this form, over all
choices, is at most \(2^\lambda\)
\cite[p.~230]{Shelah1971Ultrapowers}.
\end{hblemma}

Let \(\langle(M_i,N_i):i<2^\lambda\rangle\) list, possibly with repetitions,
all pairs of elementarily equivalent \(L_i\)-models with \(L_i\subseteq L^0\)
whose universes are cardinals \({<}\,\mu\), and let
\((\dom M_i)^\lambda=\{\bar a^i_\alpha:\alpha<2^\lambda\}\) and
\((\dom N_i)^\lambda=\{\bar b^i_\alpha:\alpha<2^\lambda\}\)
\cite[p.~230]{Shelah1971Ultrapowers}; each set has cardinality
\(\chi^\lambda\le2^\lambda\) for \(\chi=\lVert M_i\rVert\) or
\(\lVert N_i\rVert\), and the enumerations may repeat.

\begin{hblemma}{thm:p03-w03-p03-017-bookkeeping}{proved=annex:p03-w03-p03-017-bookkeeping}{The bookkeeping function \(R\)}
There is a function \(R\) on the ordinals \(\gamma<2^\lambda\) such that:
\begin{enumerate}
  \item[A)] for every \(i,\alpha<2^\lambda\) some \(\gamma<2^\lambda\) has
        \(R(\gamma)=\langle i,1,\bar a^i_\alpha\rangle\);
  \item[B)] for every \(i,\alpha<2^\lambda\) some \(\gamma<2^\lambda\) has
        \(R(\gamma)=\langle i,2,\bar b^i_\alpha\rangle\);
  \item[C)] for every \(i<2^\lambda\) and every set \(p\) of formulas with
        \(2^{\abs p}\le2^\lambda\), closed under conjunctions, of the form
        \(p=\{\phi_l(x,y_{\alpha_{l,1}},\ldots,y_{\alpha_{l,n_l}}):l<\abs p\}\)
        with \(\alpha_{l,m}<2^\lambda\) and \(\phi_l\) in \(L_i\), there are
        \(\gamma<2^\lambda\) and ordinals \(\gamma_{\alpha_{l,m}}<\gamma\)
        with \(R(\gamma)=\langle i,p\rangle\) and
        \(R(\gamma_{\alpha_{l,m}})=\langle i,1,\bar a^i_{\alpha_{l,m}}\rangle\);
        moreover, whenever \(R(\gamma)=\langle i,p\rangle\), each parameter
        function \(\bar a^i_{\alpha_{l,m}}\) of \(p\) satisfies
        \(R(\gamma')=\langle i,1,\bar a^i_{\alpha_{l,m}}\rangle\) for some
        \(\gamma'<\gamma\);
  \item[D)] for every \(A\subseteq\lambda\) some \(\gamma<2^\lambda\) has
        \(R(\gamma)=A\);
  \item[E)] for every \(\gamma\) exactly one of the forms in A)--D) occurs.
\end{enumerate}
\cite[pp.~230--231]{Shelah1971Ultrapowers}.  The variable \(y_\alpha\) in C)
stands for the parameter \(\bar a^i_\alpha/D\).
\end{hblemma}

\begin{hbtheorem}{thm:p03-w03-p03-017-induction}{proved=annex:p03-w03-p03-017-induction}{The inductive construction}
There are, for \(\gamma\le2^\lambda\), sets of functions \(F_\gamma\),
filters \(D_\gamma\) over \(\lambda\), and sets of pairs
\(H^i_\gamma\subseteq(\dom M_i)^\lambda\times(\dom N_i)^\lambda\) for
\(i<2^\lambda\), such that:
\begin{enumerate}
  \item for every \(\gamma\), \((F_\gamma,\varnothing,D_\gamma)\) is
        \((\lambda+\abs\gamma)\)-consistent, \(\abs{F_0}=2^\lambda\),
        \(D_0=\{\lambda\}\), \(\abs{F_0-F_\gamma}\le\lambda+\abs\gamma\),
        and \(F_\beta\supseteq F_\gamma\), \(D_\beta\subseteq D_\gamma\) for
        \(\beta<\gamma\);
  \item \(H^i_\beta\subseteq H^i_\gamma\) for \(\beta<\gamma\), and
        \(\abs{\bigcup_{i<2^\lambda}H^i_\gamma}\le\abs\gamma\);
  \item if \((\bar a_1,\bar b_1),\ldots,(\bar a_n,\bar b_n)\in H^i_\gamma\)
        and \(\phi(x_1,\ldots,x_n)\) is in \(L_i\), then
        \[
          \{k<\lambda:M_i\models\phi[\bar a_1[k],\ldots,\bar a_n[k]]
          \Leftrightarrow N_i\models\phi[\bar b_1[k],\ldots,\bar b_n[k]]\}
          \in D_\gamma;
        \]
  \item if \(\bar a_1,\ldots,\bar a_n\in\mathrm{Dom}\,H^i_\gamma\) and
        \(\phi\) is in \(L_i\), then one of the sets
        \(\{k<\lambda:M_i\models\phi[\bar a_1[k],\ldots,\bar a_n[k]]\}\) and
        \(\{k<\lambda:M_i\models\neg\phi[\bar a_1[k],\ldots,\bar a_n[k]]\}\)
        belongs to \(D_\gamma\);
  \item if \(R(\gamma)=\langle i,1,\bar a^i_\alpha\rangle\), then
        \(\bar a^i_\alpha\in\mathrm{Dom}\,H^i_{\gamma+1}\);
  \item if \(R(\gamma)=\langle i,2,\bar b^i_\alpha\rangle\), then
        \(\bar b^i_\alpha\in\mathrm{Range}\,H^i_{\gamma+1}\);
  \item if \(R(\gamma)=\langle i,p\rangle\) and, for every
        \(\phi(x,y_{\alpha_1},\ldots,y_{\alpha_r})\in p\),
        \(\{k<\lambda:M_i\models(\exists x)\phi(x,\bar a^i_{\alpha_1}[k],
        \ldots,\bar a^i_{\alpha_r}[k])\}\in D_\gamma\), then there is
        \(\bar a\in(\dom M_i)^\lambda\) such that for every such \(\phi\in p\)
        \(\{k<\lambda:M_i\models\phi[\bar a[k],\bar a^i_{\alpha_1}[k],\ldots,
        \bar a^i_{\alpha_r}[k]]\}\in D_{\gamma+1}\);
  \item if \(R(\gamma)=A\subseteq\lambda\), then \(A\in D_{\gamma+1}\) or
        \(\lambda-A\in D_{\gamma+1}\).
\end{enumerate}
\cite[p.~231, conditions 1)--8); p.~232, Cases I--VI]{Shelah1971Ultrapowers}.
\end{hbtheorem}

\begin{hbtheorem}{thm:p03-w03-p03-017-shelah-proof}{proved=annex:p03-w03-p03-017-shelah-proof}{Parts~1 and~2 of Shelah's theorem}
Let \(D=D_{2^\lambda}\) from \cref{thm:p03-w03-p03-017-induction}.  Then
\(D\) is an ultrafilter over \(\lambda\), and:
\begin{enumerate}
  \item if \(M\), \(N\) are elementarily equivalent models of power
        \({<}\,\mu\), then \(M^\lambda/D\cong N^\lambda/D\);
  \item if \(M\) is a model of power \({<}\,\mu\) and \(2^\kappa\le2^\lambda\)
        for an infinite cardinal \(\kappa\), then \(M^\lambda/D\) is
        \(\kappa^+\)-saturated.
\end{enumerate}
These are parts~1 and~2 of \cref{thm:p03-w03-p03-018-shelah}
\cite[pp.~231--232]{Shelah1971Ultrapowers}.  No assumption on \(2^\lambda\)
beyond \(\mathsf{ZFC}\) is used.
\end{hbtheorem}

\begin{hbtheorem}{thm:p03-w03-p03-017-keisler-shelah}{proved=annex:p03-w03-p03-017-keisler-shelah}{Keisler--Shelah, both directions}
Assume the Axiom of Choice.  For structures \(M\) and \(N\) of one
signature, \(M\equiv N\) if and only if there are nonempty sets \(I\), \(J\)
and ultrafilters \(U\) on \(I\) and \(V\) on \(J\) with
\(M^I/U\cong N^J/V\).  For the forward direction the construction gives
\(I=J=\lambda\) and \(U=V=D\) for any infinite \(\lambda\) with
\(\lVert M\rVert,\lVert N\rVert<\mu\), for instance
\(\lambda=2^{\lVert M\rVert+\lVert N\rVert+\aleph_0}\); nothing is asserted
about \(D\) beyond \cref{thm:p03-w03-p03-017-shelah-proof}.  This proves
\cref{thm:p03-keisler-shelah}.
\end{hbtheorem}

\begin{hbexample}{def:p03-w03-p03-017-lemma1-example}{Lemma~1 at \(\lambda=\aleph_0\)}
For \(\lambda=\aleph_0\) one has \(\mu=\aleph_0\), and Lemma~1 gives
\(2^{\aleph_0}\) functions \(f_B:\omega\to\omega\), one for each
\(B\subseteq\omega\), such that any finitely many distinct ones take any
prescribed values at a common argument.  The proof lists all pairs
\((A_k,h_k)\), \(k<\omega\), of a finite \(A\subseteq\omega\) and a function
\(h\) from a finite family of subsets of \(A\) into \(\omega\), and sets
\(f_B(k)=h_k(B\cap A_k)\) when defined and \(0\) otherwise.  Take \(B_1\) the
even numbers, \(B_2\) the multiples of \(3\), and the targets \(5\) and
\(7\).  The point \(2\) lies in \(B_1\) and not in \(B_2\); with
\(A=\{2\}\) the traces are \(B_1\cap A=\{2\}\) and \(B_2\cap A=\varnothing\),
and \(h(\{2\})=5\), \(h(\varnothing)=7\).  At the index \(k\) where
\((A_k,h_k)=(A,h)\) one has \(f_{B_1}(k)=5\) and \(f_{B_2}(k)=7\).  Since
\(\mu=\aleph_0\), part~1 of the theorem then covers only finite models,
as noted in \cref{def:p03-w03-p03-018-worked}.
\end{hbexample}

\begin{hbwarning}{def:p03-w03-p03-017-shelah-reading}{How the source is read here}
\begin{itemize}
  \item The source calls \(H^i_\gamma\) a function.  Lemma~5 does not
        exclude that the element it produces already occurs in an earlier
        pair, so \(H^i_\gamma\) is read as a set of pairs, and condition~3
        is stated for pairs.  The map induced on classes is still a well
        defined bijection by condition~3 applied to \(x_1=x_2\).
  \item In Lemma~3 the source indexes the \((n+1)\)-sequences of values
        below \(\chi^*\) by ordinals \(\beta<\chi^*\); when \(\chi^*\) is
        finite there are more such sequences than \(\chi^*\).  The proof
        uses only that their number is below \(\mu\).
  \item In Lemma~5 the source chooses \(f\in F_1\); the hypothesis
        \(\abs F>\kappa\), stated above, makes \(F_1\) nonempty, and it
        holds at every stage of the construction.
  \item The fixed language \(L^0\) is taken with \(\lambda\) constant
        symbols in addition to the source's relation and function symbols
        of each positive arity.
  \item Saturation in the handbook concerns types in finitely many
        variables (\cref{thm:saturated-extension}); condition~7 realizes
        one-variable types, and the passage to finitely many variables is
        written out in the proof.
\end{itemize}
\end{hbwarning}

\begin{proofboundary}{def:p03-w03-p03-017-shelah-boundary}{Part~3 and goodness}
Part~3 of \cref{thm:p03-w03-p03-018-shelah} (isomorphism of elementarily
equivalent ultraproducts \(\prod_{k<\lambda}M_k/D\),
\(\prod_{k<\lambda}N_k/D\) of models of cardinality at most \(\chi<\mu\)) is
stated in the source, whose proof of the main theorem omits it
\cite[p.~230, Remark~2]{Shelah1971Ultrapowers}; it is not proved here.  The
source also remarks that \(D\) is \(\lambda^+\)-good, by part~2 and an
earlier result of Keisler, or directly as in Kunen
\cite[p.~230, Remark~3]{Shelah1971Ultrapowers}; that remark is not proved
here either.  The source's closing remark that the lemmas, except Lemma~2B,
extend to fewer than \(\chi^0\) equations \(f(k)=g(k)\) in Definition~1
\cite[p.~232]{Shelah1971Ultrapowers} is likewise only reported.
\end{proofboundary}

\subsection{Keisler's 1961 characterizations of elementary classes}
\label{sec:p03-w03-p03-017-keisler1961}

Keisler's conditional isomorphism theorem, its set-theoretic lemma, its
ultrapower case, its existential form under the generalized continuum
hypothesis and the ultrafilter-first forms announced in its appendix are
stated and proved, at the article's own hypotheses and in its notation, in
\cref{sec:p03-w03-p03-018-keisler}
(\cref{thm:p03-w03-p03-018-keisler-disjoint,thm:p03-w03-p03-018-keisler-products,thm:p03-w03-p03-018-keisler-powers,thm:p03-w03-p03-018-keisler-gch,thm:p03-w03-p03-018-keisler-appendix}).
This subsection adds the rest of the article's Section~2: the
characterizations of elementary classes that it derives from those theorems.
The notation is that of \cref{def:p03-w03-p03-018-keisler-notation}, where
the similarity type is left unnamed.

The article states in its first footnote that it is identical with Keisler's
doctoral dissertation, submitted at the University of California in the
spring of 1961 \cite[p.~477, footnote~*]{Keisler1961ElementaryClasses}.
Shelah's bibliography lists the dissertation itself as its item~6 and gives
the journal article only as a further reference
\cite[p.~232, reference~6]{Shelah1971Ultrapowers}.  The account here reads
the journal article, pages 477--495; the dissertation is a separate document
that was not read, and nothing here is attributed to it.

\begin{hbconvention}{def:p03-w03-p03-017-keisler-classes}{Classes of structures in Keisler's article}
All classes consist of structures of one fixed relational type.  For a
class \(K\), \(K^*\) is the set of sentences true in every member of \(K\),
and \(\bar K\) is the class of structures of the type that are not in \(K\).
A class is in \(EC\) if it is the class of models of one sentence, in
\(EC_\Delta\) if it is the class of models of a set of sentences, in
\(EC_\Sigma\) if it is a union of classes in \(EC\), equivalently if its
complement is in \(EC_\Delta\), and in
\(EC_{\Delta\Sigma}\) if it is a union of classes in \(EC_\Delta\);
equivalently, \(K\in EC_{\Delta\Sigma}\) exactly when \(K\) is closed under
elementary equivalence, and such classes are called \emph{elementarily
closed}.  The \emph{elementarily closed class generated by} \(M\) is the
class of structures elementarily equivalent to some member of \(M\)
\cite[pp.~481--482]{Keisler1961ElementaryClasses}.  \(\mathrm{Prod}(K)\) is
the class of structures isomorphic to an ultraproduct of a system of members
of \(K\), and \(\mathrm{Pow}(K)\) the class of structures isomorphic to an
ultrapower of a member of \(K\) \cite[p.~483]{Keisler1961ElementaryClasses}.
\end{hbconvention}

The article states Theorems~1.1--1.6 without proof, referring to Frayne,
Morel and Scott and to Kochen \cite[p.~483]{Keisler1961ElementaryClasses}.
Among them, Theorem~1.2 says that \(K\in EC_{\Delta\Sigma}\) implies
\(K=\mathrm{Pow}(K)\); Theorem~1.3 that \(K\in EC_\Delta\) if and only if
\(K\in EC_{\Delta\Sigma}\) and \(K=\mathrm{Prod}(K)\); and Theorem~1.4 that
\(K\in EC\) if and only if \(K\in EC_{\Delta\Sigma}\),
\(K=\mathrm{Prod}(K)\) and \(\bar K=\mathrm{Prod}(\bar K)\).  After
Theorem~2.3 the article assumes the generalized continuum hypothesis for the
rest of Section~2 \cite[p.~487]{Keisler1961ElementaryClasses} and states:
\begin{itemize}
  \item Theorem~2.4, the existential form
        (\cref{thm:p03-w03-p03-018-keisler-gch}).
  \item Theorem~2.5: \(K^*\cup M^*\) is consistent if and only if
        \(\mathrm{Prod}(K)\cap\mathrm{Prod}(M)\neq\varnothing\).  Its proof
        applies the ultraproduct theorem
        \cref{thm:p03-w03-p03-018-keisler-products} to a system indexed by a
        set of cardinality \(\alpha\) built from members of \(K\) and \(M\)
        chosen by the Axiom of Choice
        \cite[pp.~487--488]{Keisler1961ElementaryClasses}.
  \item Theorem~2.6 and Corollary~2.7: \(K\in EC_{\Delta\Sigma}\) if and
        only if \(K=\mathrm{Pow}(K)\) and \(\bar K=\mathrm{Pow}(\bar K)\),
        if and only if
        \(K=\{\mathfrak A:\mathrm{Pow}(\mathfrak A)\cap\mathrm{Pow}(M)\neq
        \varnothing\}\) for some \(M\subseteq K\); and the elementarily
        closed class generated by \(M\) is
        \(\{\mathfrak A:\mathrm{Pow}(\mathfrak A)\cap\mathrm{Pow}(M)\neq
        \varnothing\}\) \cite[p.~488]{Keisler1961ElementaryClasses}.
  \item Theorem~2.8, Corollaries~2.9 and~2.10 and Theorem~2.11, the
        corresponding characterizations of \(EC_\Delta\), \(EC_\Sigma\) and
        \(EC\) by \(\mathrm{Prod}\) and \(\mathrm{Pow}\).  Theorem~2.8 and
        its corollaries use Theorem~2.5.  Theorem~2.11 states that
        \(K\in EC\) if and only if \(K=\mathrm{Prod}(K)\) and
        \(\bar K=\mathrm{Prod}(\bar K)\), and its proof uses only
        Theorem~2.6 and Theorem~1.4 \cite[p.~488]{Keisler1961ElementaryClasses}.
  \item Theorems~2.12--2.14, separation principles, of which the last
        yields Craig's interpolation theorem; the article remarks that the use
        of the generalized continuum hypothesis there does no harm, because
        the syntactic form of Craig's theorem is a statement about natural
        numbers, and such statements provable with the generalized continuum
        hypothesis and Choice are provable without them
        \cite[pp.~488--490]{Keisler1961ElementaryClasses}.
\end{itemize}
The article closes Section~2 by asking whether any of 2.4--2.14 can be
proved without the generalized continuum hypothesis
\cite[p.~490]{Keisler1961ElementaryClasses}.  Its appendix also states,
under that hypothesis and without proofs, Horn characterizations: with
\(HC_\Delta\) the classes axiomatized by sets of Horn sentences and
\(\mathrm{RProd}(K)\) the class of structures isomorphic to reduced products
of members of \(K\), \(K\in HC_\Delta\) if and only if
\(\mathrm{RProd}(K)=K\) and \(\mathrm{Pow}(K)=K\) (Theorem~A.10); the
article attributes to Chang the direction that \(K\in HC_\Delta\) implies
\(\mathrm{RProd}(K)=K\) \cite[p.~493]{Keisler1961ElementaryClasses}.

\subsection{Earlier approaches and consequences in Shelah's account}
\label{sec:p03-w03-p03-017-shelah-context}

The first two pages of \textcite{Shelah1971Ultrapowers} place the theorem
among earlier work.  In those passages the source uses \(\mu\) also for an
arbitrary cardinal, before its notation paragraph fixes
\(\mu=\min\{\mu:\lambda^\mu>\lambda\}\); below, a \(\mu\) quoted from those
passages is an arbitrary infinite cardinal, and this is said each time.

\paragraph{Keisler's theorem and its decomposition.}  Shelah reports that
Keisler's item~6 proves, assuming \(\lambda^+=2^\lambda\), an ultrafilter
\(D\) over \(\lambda\) with \(M^\lambda/D\cong N^\lambda/D\) whenever
\(M\equiv N\), \(\lVert M\rVert,\lVert N\rVert\le\lambda^+\), and the
language has cardinality \({\le}\,\lambda\)
\cite[p.~224]{Shelah1971Ultrapowers}.  The article proves this with two
ultrafilters chosen for the given pair
(\cref{thm:p03-w03-p03-018-keisler-powers}); the single ultrafilter is the
announced form, item~2 of \cref{thm:p03-w03-p03-018-keisler-appendix}, as
\cref{sec:p03-w03-p03-018-keisler} explains.  Shelah then
reports that, by Keisler's item~8, the result splits into three stages
\cite[p.~224]{Shelah1971Ultrapowers}:
\begin{enumerate}
  \item if \(\lambda^+=2^\lambda\), there is a \(\lambda^+\)-good
        ultrafilter over \(\lambda\);
  \item if \(D\) is a \(\lambda^+\)-good ultrafilter over a set \(I\) and
        \(M\) has a language of cardinality \({\le}\,\lambda\), then
        \(M^I/D\) is \(\lambda^+\)-saturated;
  \item any two elementarily equivalent \(\mu\)-saturated models of
        cardinality \(\mu\) are isomorphic, for an arbitrary infinite
        cardinal \(\mu\).
\end{enumerate}
The source refers for these to Keisler's items~7 and~8 and to Morley and
Vaught (its item~15).  It does not define \(\lambda^+\)-good on these pages,
and stages~1 and~2 are reported here as Shelah's account only.  Stage~3 and
the way the stages combine are proved below.

\begin{hbtheorem}{thm:p03-w03-p03-017-saturated-unique}{proved=annex:p03-w03-p03-017-saturated-unique}{Uniqueness of saturated models of one power}
Assume the Axiom of Choice.  Let \(\mu\) be an infinite cardinal, here
arbitrary, and let \(M\), \(N\) be elementarily equivalent structures of one
signature, each of cardinality \(\mu\) and \(\mu\)-saturated.  Then
\(M\cong N\).  This is stage~3 of the decomposition reported in
\cite[p.~224]{Shelah1971Ultrapowers}.
\end{hbtheorem}

\begin{hbcorollary}{thm:p03-w03-p03-017-decomposition}{proved=annex:p03-w03-p03-017-decomposition}{How the stages give Keisler's theorem}
Let \(\lambda\) be an infinite cardinal with \(\lambda^+=2^\lambda\), and let
\(D\) be an ultrafilter over \(\lambda\) such that \(M^\lambda/D\) is
\(\lambda^+\)-saturated for every model \(M\) whose language has cardinality
\({\le}\,\lambda\).  If \(M\equiv N\), \(\lVert M\rVert,\lVert N\rVert\le
\lambda^+\) and the language has cardinality \({\le}\,\lambda\), then
\(M^\lambda/D\cong N^\lambda/D\).
\end{hbcorollary}

By stage~2, a \(\lambda^+\)-good ultrafilter over \(\lambda\) has the
property assumed of \(D\), and by stage~1 one exists when
\(\lambda^+=2^\lambda\); so the corollary turns the reported stages~1 and~2
into Keisler's theorem as Shelah states it.

\paragraph{Ultralimits and Boolean ultrapowers.}  Shelah reports two further
approaches \cite[pp.~224--225]{Shelah1971Ultrapowers}.  Kochen (Shelah's
item~11; also Keisler's item~10, \S5) replaced ultrapowers by ultralimits,
which keep most of the useful properties of ultrapowers, and proved that any
two elementarily equivalent models of cardinalities \({\le}\,\beth_\alpha\)
have isomorphic ultralimits of cardinality \(\beth_{\alpha+\omega}\).
Mansfield (item~13) generalized ultrapowers to Boolean ultrapowers and proved
a parallel isomorphism theorem for them.  Keisler's 1961 introduction
likewise mentions that two structures are elementarily equivalent exactly
when they have isomorphic limit ultrapowers, a result that needs no
continuum hypothesis but uses a more complex operation
\cite[p.~478]{Keisler1961ElementaryClasses}.  Keisler's footnote~5 adds
that Kochen announced Theorem~2.4 at the end of 1960 as a consequence of a
stronger statement, which Keisler shows false by a counterexample, so that
the proof Kochen had in mind does not seem correct
\cite[p.~487, footnote~5]{Keisler1961ElementaryClasses}.  None of these
results is proved in this handbook.

\paragraph{Kunen, Silver and Rucker, and an open question.}  Kunen (item~12)
eliminated the generalized continuum hypothesis from the existence of good
ultrafilters, and the construction of
\cref{sec:p03-w03-p03-017-shelah} generalizes his proof
\cite[p.~225]{Shelah1971Ultrapowers}.  Silver and, independently, Rucker
proved that it is consistent with \(\mathsf{ZFC}+\aleph_1<2^{\aleph_0}\)
that some ultrafilter \(D\) over \(\omega\) makes \(M^\omega/D\) saturated
for every countable model \(M\) with countable language; Shelah adds that
this follows readily from Martin's axiom
\cite[p.~225]{Shelah1971Ultrapowers}.  These results are reported, not
proved here.  Shelah then states as open whether, for all \(M\equiv N\) with
\(\lVert M\rVert,\lVert N\rVert\le\mu\) and \(\abs{L(N)}\le\mu\), for an
arbitrary infinite cardinal \(\mu\), there is an ultrafilter \(D\) over
\(\mu\) with \(M^\mu/D\cong N^\mu/D\), and suggests that this may be
independent of \(\mathsf{ZFC}\) \cite[p.~225]{Shelah1971Ultrapowers}.  The
question asks for the index set equal to the size of the models; the theorem
of \cref{sec:p03-w03-p03-018-shelah-ultrapowers} instead needs models of
power below the cardinal \(\mu\) attached to the index set \(\lambda\), and
\(\mu\le\lambda\) by \cref{thm:p03-w03-p03-018-mu}.  Its status after 1971
is not recorded here.

\paragraph{Eliminating the generalized continuum hypothesis.}  Shelah states
that part~1 of the theorem removes the generalized continuum hypothesis from
some of the theorems Keisler used in his item~6, in particular those
characterizing elementary classes, and that the technique of the paper
removes it from Keisler's theorem (his item~5) that a sentence is preserved
under reduced products if and only if it is equivalent to a Horn sentence;
for the latter he notes that Galvin (item~3) had already done so by a
set-theoretic argument and Mansfield (item~14) by Boolean ultrapowers
\cite[p.~225]{Shelah1971Ultrapowers}.  The source gives no proof of either
consequence.  The part that follows from part~1 alone is proved here.

Theorem~2.4 itself, with the transitivity that Keisler left open, holds
without any cardinal hypothesis: this is item~3 of
\cref{thm:p03-w03-p03-018-keisler-compare}, which rests on
\cref{thm:p03-w03-p03-017-keisler-shelah}.  The characterizations that
depend only on Theorem~2.4 follow.

\begin{hbtheorem}{thm:p03-w03-p03-017-gch-free-closed}{proved=annex:p03-w03-p03-017-gch-free-closed}{Keisler's Theorems~2.6, 2.11 and~2.12 and Corollary~2.7 without the continuum hypothesis}
Assume \(\mathsf{ZFC}\) only, and let \(K\), \(M\) be classes of structures
of one relational similarity type.
\begin{enumerate}
  \item The following are equivalent: (i) \(K\in EC_{\Delta\Sigma}\);
        (ii) \(K=\mathrm{Pow}(K)\) and \(\bar K=\mathrm{Pow}(\bar K)\);
        (iii) for some \(M\subseteq K\),
        \(K=\{\mathfrak A:\mathrm{Pow}(\mathfrak A)\cap\mathrm{Pow}(M)\neq
        \varnothing\}\).
  \item The elementarily closed class generated by \(M\) equals
        \(\{\mathfrak A:\mathrm{Pow}(\mathfrak A)\cap\mathrm{Pow}(M)\neq
        \varnothing\}\).
  \item If \(K=\mathrm{Pow}(K)\), \(M=\mathrm{Pow}(M)\) and
        \(K\cap M=\varnothing\), then the elementarily closed class
        generated by \(K\) contains \(K\) and is disjoint from \(M\).
  \item \(K\in EC\) if and only if \(K=\mathrm{Prod}(K)\) and
        \(\bar K=\mathrm{Prod}(\bar K)\).
\end{enumerate}
Item~4 is the characterization of elementary classes.  The proof writes out
the parts of Keisler's Theorems~1.2--1.4 that it needs, which the article
states without proof \cite[p.~483]{Keisler1961ElementaryClasses}, and uses
no cardinal hypothesis.
\end{hbtheorem}

\begin{proofboundary}{def:p03-w03-p03-017-context-boundary}{What is reported and not proved}
\begin{itemize}
  \item Stages~1 and~2 of the decomposition (existence of
        \(\lambda^+\)-good ultrafilters under \(\lambda^+=2^\lambda\), and
        \(\lambda^+\)-saturation of ultrapowers by them) are Shelah's
        account of Keisler's items~7 and~8, which are not among the held
        sources; the notion of a good ultrafilter is not defined here.
  \item The results of Kochen, Mansfield, Kunen, Silver and Rucker are
        reported as the source states them; their proofs are not given
        here.
  \item Keisler's Theorems~2.5, 2.8, 2.13 and~2.14 and Corollaries~2.9
        and~2.10 rest on Theorem~2.5,
        whose proof applies Theorem~2.2 to ultraproducts of families of
        structures, not to ultrapowers of one structure.  Part~1 of Shelah's
        theorem concerns ultrapowers, and the ultraproduct form, part~3, is
        not proved in the source.  The versions of those theorems without
        the continuum hypothesis are therefore not established here.
  \item The removal of the continuum hypothesis from the Horn preservation
        theorem is reported only.  The direction proved in this handbook
        without any cardinal hypothesis is that Horn formulas are preserved
        by reduced products (\cref{thm:p03-w03-p03-017-horn-preservation}).
  \item Shelah's p.~225 recommends the account of \L o\'s as item~3 of his
        list; in that list item~3 is Galvin's paper and \L o\'s's paper is
        item~4 \cite[pp.~225 and~232]{Shelah1971Ultrapowers}.  Keisler's
        proof of Theorem~2.2 cites Theorem~1.4 for the reducts of the
        expanded ultraproducts; the statement used there is Theorem~1.5
        \cite[pp.~484 and~486]{Keisler1961ElementaryClasses}.
\end{itemize}
\end{proofboundary}

\paragraph{Relationships.}  The statement, its hypotheses and the comparison
with Keisler's conditional theorem are in
\cref{sec:p03-w03-p03-018-shelah-ultrapowers}; the converse direction there
is item~2 of \cref{thm:p03-w03-p03-018-comparison}, and the converse step of
\cref{thm:p03-w03-p03-017-keisler-shelah} is the same argument.  Ultrapowers
and \L o\'s's theorem are \cref{sec:types-ultraproducts-05}; saturation is
the notion preceding \cref{thm:saturated-extension}.  K\"onig's theorem
(\cref{thm:p06-konig-cardinal}) enters only the bookkeeping lemma.  The
choice assumptions are those of \cref{def:p03-ultrafilter-choice-boundary}.
Keisler's \cref{thm:p03-w03-p03-018-keisler-products} and Shelah's
\cref{thm:p03-w03-p03-017-shelah-proof} are both back-and-forth
constructions of the same shape; the first needs
\(\varkappa(2^\alpha)=\alpha^+\) to enumerate the products in order type
\(\alpha^+\), the second replaces that enumeration by the independent family
of Lemma~1 and fixes the ultrafilter before the models.  The countable case of
\cref{thm:p03-w03-p03-017-saturated-unique} is
\cref{thm:p03-vaught-saturated-uniqueness}.  Compactness enters the
elementary-class characterization through
\cref{thm:p03-compactness-ultraproduct}.

\paragraph{Sources.}  Horn formulas, direct and reduced products and the
preservation theorem: S.~Burris and H.~P.~Sankappanavar, \emph{A Course in
Universal Algebra}, Millennium Edition (2012 update), Chapter~V, \S1
(Definition~1.4 on p.~192, direct products and Exercises~29--30 on p.~204)
and \S2 (Definitions~2.1--2.5 and Lemma~2.6 on pp.~205--207, Theorem~2.7 on
pp.~207--210) \cite{BurrisSankappanavar2012}.  The construction: S.~Shelah,
\emph{Every two elementarily equivalent models have isomorphic
ultrapowers}, Israel Journal of Mathematics 10 (1971), pages 224--233
\cite{Shelah1971Ultrapowers}: notation p.~225, Definition~1 and Lemmas~1--2
pp.~226--227, Lemmas~3--4 pp.~227--229, Lemma~5 pp.~229--230, the proof of
the main theorem pp.~230--232; the earlier approaches, the open question and
the consequences pp.~224--225; the bibliography pp.~232--233.  Keisler's
construction and characterizations: H.~J.~Keisler, \emph{Ultraproducts and
elementary classes}, Koninklijke Nederlandse Akademie van Wetenschappen,
Proceedings Series~A 64 (Indagationes Mathematicae 23) (1961), pages
477--495 \cite{Keisler1961ElementaryClasses}: notation pp.~479--483,
Lemma~2.1 and Theorem~2.2 pp.~484--486, Theorems~2.3--2.14 pp.~486--490, the
appendix pp.~490--493.  Both presentations follow the published journal
pages.
